% Part: computability
% Chapter: recursive-functions
% Section: composition

\documentclass[../../../include/open-logic-section]{subfiles}

\begin{document}

\olfileid{cmp}{rec}{com}
\olsection{సంయుక్తం}

$f$, $g$ సహజ సంఖ్యల రెండు ఏకస్థానిక ప్రమేయాలైతే, వాటి
సంయుక్తాన్ని తీసుకోవచ్చు: $h(x)=f(g(x))$. అప్పుడు కొత్త ప్రమేయం~$h(x)$,
$f$,~$g$ల నుంచి \emph{సంయుక్తం} ద్వారా నిర్వచితమవుతుంది. దీన్ని ఒకటి
కంటే ఎక్కువ ఆర్గుమెంటులు గల ప్రమేయాలకు సాధారణీకరించాలని కోరుకుంటాం.

దీనికి ఒక పద్ధతి ఇది: $f$ ఒక $k$-స్థానిక ప్రమేయం, $g_0$, \dots,
$g_{k-1}$ అన్నీ $n$-స్థానికాలైన $k$ ప్రమేయాలు అనుకుందాం. అప్పుడు
కొత్త $n$-స్థానిక ప్రమేయం $h$ను కింది విధంగా నిర్వచించవచ్చు:
\[
h(x_0, \dots, x_{n-1}) =
f(g_0(x_0, \dots, x_{n-1}), \dots, g_{k-1}(x_0, \dots, x_{n-1}))
\]
$f$, అన్ని $g_i$లు గణనీయమైతే, $h$ కూడా గణనీయమే:
$h(x_0,\dots,x_{n-1})$ను గణించడానికి, మొదట ప్రతి $i=0$, \dots,
~$k-1$కూ $y_i=g_i(x_0,\dots,x_{n-1})$ విలువలను గణించాలి. తరువాత ఈ
విలువలను $f$కు ఇచ్చి $h(x_0,\dots,x_{n-1})=f(y_0,\dots,y_{k-1})$ను
గణించాలి.
\sourcecorrection{OLTECMPREC-002}{ఆంగ్ల మూలంలోని గణన సూచన $h$కు $n$ ఆర్గుమెంటులు ఉండాల్సిన చోట $x_{k-1}$తో ముగిసింది; నిర్వచక ప్రదర్శనకు అనుగుణంగా $x_{n-1}$గా సరిచేశాం.}

ఇప్పటికే ఉన్న ప్రమేయాలను ఉపయోగించి కొత్త ప్రమేయాన్ని గణించేటప్పుడు
ఏమి జరుగుతుందో చెప్పడానికి ఇది మరీ నిర్బంధకమైన లక్షణీకరణలా
కనిపించవచ్చు. ఉదాహరణకు, $\Add$ యొక్క ఆదిమ పునరావృత్త నిర్వచనంలో
$g(x,y,z)=\Succ(z)$ను వాడినప్పుడు చేసినట్లుగా, కొన్నిసార్లు ప్రమేయపు
ఆర్గుమెంటులన్నిటినీ వాడం. కింది ప్రమేయాలను వాడేందుకు మనకు అనుమతి
ఉందనుకుందాం:
\[
\Proj{n}{i}(x_0, \dots, x_{n-1}) = x_i
\]
$\Proj{n}{i}$లను \emph{ప్రక్షేప ప్రమేయాలు} అంటారు:
\sourcecorrection{OLTECMPREC-005}{ఆంగ్ల మూలం ఈ వాక్యంలో ప్రక్షేప ప్రమేయాన్ని $\Proj{k}{i}$గా సూచించింది; వెంటనే ముందున్న నిర్వచనం, తరువాతి స్థానసంఖ్య వివరణలకు అనుగుణంగా $\Proj{n}{i}$గా సరిచేశాం.}
$\Proj{n}{i}$ ఒక $n$-స్థానిక ప్రమేయం. అప్పుడు $g$ను
\[
g(x, y, z) = \Succ(\Proj{3}{2}(x, y, z)).
\]
అని నిర్వచించవచ్చు. ఇక్కడ $f$ పాత్రను $1$-స్థానిక ప్రమేయం $\Succ$
పోషిస్తుంది కనుక $k=1$. $g_0$ పాత్రను పోషించే $3$-స్థానిక ప్రమేయం
$\Proj{3}{2}$ ఒక్కటే ఉంది. ఫలితం మూడవ ఆర్గుమెంటు ఉత్తరగామిని ఇచ్చే
$3$-స్థానిక ప్రమేయం.

ఆర్గుమెంటుల క్రమాన్ని మార్చడం లేదా వాటిని ఒకటిగా గుర్తించడం ద్వారా
కొత్త ప్రమేయాలను నిర్వచించడానికీ ప్రక్షేప ప్రమేయాలు వీలు కల్పిస్తాయి.
ఉదాహరణకు, $h(x)=\Add(x,x)$ ప్రమేయాన్ని
\[
h(x_0) = \Add(\Proj{1}{0}(x_0),\Proj{1}{0}(x_0)).
\]
అని నిర్వచించవచ్చు. ఇక్కడ $k=2$, $n=1$; $f(y_0,y_1)$ పాత్రను
$\Add$ పోషిస్తుంది; $g_0(x_0)$, $g_1(x_0)$ పాత్రలు రెండింటినీ
ఏకస్థానిక ప్రక్షేప ప్రమేయం (అంటే తాదాత్మ్య ప్రమేయం)
$\Proj{1}{0}(x_0)$ పోషిస్తుంది.

$f(y_0,y_1)$ ఇప్పటికే మనకు ఉన్న ప్రమేయమైతే, $h(x_0,x_1)=f(x_1,x_0)$
ప్రమేయాన్ని
\[
h(x_0, x_1) = f(\Proj{2}{1}(x_0, x_1),\Proj{2}{0}(x_0, x_1)).
\]
అని నిర్వచించవచ్చు. ఇక్కడ $k=2$, $n=2$; $g_0$, $g_1$ పాత్రలను
వరుసగా $\Proj{2}{1}$,~$\Proj{2}{0}$ పోషిస్తాయి.

$g_0$, \dots,~$g_{k-1}$ అన్నిటికీ ఒకే స్థానసంఖ్య~$n$ ఉండాలన్న
నిబంధన గురించి కూడా సందేహం రావచ్చు. (ప్రమేయపు \emph{స్థానసంఖ్య (అరిటీ)} అంటే
దాని ఆర్గుమెంటుల సంఖ్య; $n$-స్థానిక ప్రమేయానికి స్థానసంఖ్య~$n$.) అయితే
ప్రక్షేప ప్రమేయాలను చేర్చడం కావలసిన వెసులుబాటును ఇస్తుంది. ఉదాహరణకు,
$f$,~$g$ $3$-స్థానిక ప్రమేయాలు, $h$ కింది విధంగా నిర్వచించిన
$2$-స్థానిక ప్రమేయం అనుకుందాం:
\[
h(x,y) = f(x,g(x,x,y),y).
\]
$h$ నిర్వచనాన్ని ప్రక్షేప ప్రమేయాలతో ఇలా తిరిగి రాయవచ్చు:
\[
h(x,y) = f(\Proj{2}{0}(x,y), g(\Proj{2}{0}(x,y), \Proj{2}{0}(x,y),
\Proj{2}{1}(x,y)), \Proj{2}{1}(x,y)).
\]
అప్పుడు $h$ అనేది $f$, $\Proj{2}{0}$, $l$, $\Proj{2}{1}$ల
సంయుక్తం; ఇక్కడ
\[
l(x,y) = g(\Proj{2}{0}(x,y),\Proj{2}{0}(x,y),\Proj{2}{1}(x,y)),
\]
అంటే $l$, $g$, $\Proj{2}{0}$, $\Proj{2}{0}$,
~$\Proj{2}{1}$ల సంయుక్తం.

\end{document}
